\section{Results}
\label{sec:results}
\subsection{Baseline, Isolation, and Loop Rejection}
All six saved matrix runs report 31 research channels, zero build errors, and
zero build warnings. The recorded sample counts match their respective
simulation durations at the specified output interval
(Table~\ref{tab:scenarios}). This is evidence of successful case execution and
recorder extraction, not proof that the parameter set is physically calibrated.

In S1, the downstream final RMS voltages range from approximately 382.0 to
386.97 V at N03--N06. The normally open tie currents are negligible, and no
retained below-threshold interval is reported. S2 opens the three incident
branches of N03 at 5 s. N03--N06 then remain below the analysis threshold to
the 6-s run end. N04 first crosses below 200 V at 5.0165 s, giving an observed
duration of 0.9835 s. This is only a lower bound on the unresolved
interruption, not a recovery time.

S4 rejects the single loop-forming closure and produces no state-changing
switching events. Its final voltage profile remains consistent with the
baseline. This establishes the intended behavior for that particular cycle
test. It does not establish atomic rejection for a mixed command containing
both acceptable and unacceptable operations.

\subsection{Tie Restoration and Imposed Delay}
In S3, tie closure at 6 s restores N04, N05, and N06 while N03 stays isolated.
The respective final voltages are 380.51, 387.37, and 383.43 V
(Table~\ref{tab:voltages}). The final tie currents are 32.54 A on N02--N05
and 13.45 A on N04--N06. These current magnitudes provide electrical evidence
that the alternate branches carry load after closure, beyond a graphical
indication that an edge is closed. They do not establish a thermal margin
against measured cable ratings.

\begin{table}[tb]
\centering\small
\caption{Saved final bus RMS voltage (V). The small N03 residual after
restoration is not usable service; finite off-state paths remain in the model.}
\label{tab:voltages}
\begin{tabular}{@{}lrrr@{}}
\toprule
Bus & S1 baseline & S3 restoration & S6 physical fault \\
\midrule
N03 & 386.97 & 0.02 & 0.02 \\
N04 & 384.83 & 380.51 & 380.51 \\
N05 & 383.51 & 387.37 & 387.37 \\
N06 & 382.03 & 383.43 & 383.43
\\
\bottomrule
\end{tabular}
\end{table}

N04 and N06 remain below 200 V for approximately 0.9984 s in S3, and N05 for
0.99805 s. In S5, changing the tie-closure time from 6 to 45 s increases each
of these saved interval durations by 39 s; their final voltage values remain
the same to the displayed precision. The difference is the direct consequence
of the imposed restoration schedule. It cannot be described as an SDN speedup
over a measured legacy network, because neither network-delay distribution nor
communications implementation is varied in these runs.

\begin{table}[tb]
\centering\small
\caption{Completed downstream intervals below 200 V, from saved full-rate
metrics. Displayed decimal places reflect sample-based accounting rather than
physical calibration accuracy.}
\label{tab:interruptions}
\begin{tabular}{@{}lrrr@{}}
\toprule
Bus & S3 duration (s) & S5 duration (s) & S6 duration (s) \\
\midrule
N04 & 0.99840 & 39.99840 & 1.19485 \\
N05 & 0.99805 & 39.99805 & 1.19455 \\
N06 & 0.99840 & 39.99840 & 1.19495
\\
\bottomrule
\end{tabular}
\end{table}

\subsection{Native Physical Fault and Restoration}
The physical-fault scenario produces a different voltage trajectory from
isolation-only switching. The fault input is active from 4.8 to 5.2 s;
N03--N06 cross below threshold around 4.820 s, before the scheduled breaker
opening at 5 s. The maximum absolute fault phase currents are 1.552, 1.554,
and 1.552 kA for A, B, and C, respectively. The largest recorded phase peak
is 1.554333 kA. In contrast, inactive-fault cases exhibit approximately
1.3 mA maximum leakage rather than a high-current event.

Opening the three N03 incident branches removes the source-fed route to the
fault location. The fault logic clears at 5.2 s, and the healthy downstream
buses are reconnected through the two ties at 6 s. Their below-threshold
durations are 1.19485 s at N04, 1.19455 s at N05, and 1.19495 s at N06.
Their final voltages and tie currents agree with S3 at the reported precision.
N03 remains isolated to the 7-s run end, with only a small residual measured
voltage. Restoring healthy loads therefore does not imply restoration of the
faulted bus.

\begin{figure}[tb]
\centering\includegraphics[width=0.96\textwidth]{physical_fault.pdf}
\caption{S6 physical-fault response near the switching events. The upper two
panels use the stored 501-point voltage and RMS-current previews from a
140,001-sample-per-channel run; they are decimated trajectories, not full-rate
transient waveforms. The lower panel shows exact input schedules, not measured
breaker-contact feedback. Dotted lines mark 4.8, 5.0, 5.2, and 6.0 s.}
\label{fig:physical}
\end{figure}

Figure~\ref{fig:physical} combines the stored RMS trajectories with the exact
scenario input schedule. Figure~\ref{fig:comparison} reports the interval
durations and fault-current peaks from the stored full-rate metrics. The
roughly 15--20 ms difference between a programmed transition and an RMS
threshold crossing should not be interpreted as communication latency. It
includes the voltage measurement's response and the electrical transient, as
well as the threshold definition.

\begin{figure}[tb]
\centering\includegraphics[width=\textwidth]{comparison.pdf}
\caption{Saved full-rate summary metrics. Panel (a) uses a logarithmic duration
axis to show the intentionally imposed long wait in S5 alongside S3 and S6.
Panel (b) reports phase-current peaks for S6; no peaks are estimated from
decimated waveform previews.}
\label{fig:comparison}
\end{figure}

\subsection{Recorded Dashboard Integration Evidence}
The graphical-validation artifact dated 5 September 2026 records three
successful actions: selecting the baseline, selecting the physical fault, and
editing a controller command. All three retain the same PSCAD process
identifier and return 31 selected channels with no reported errors. The
baseline has 20,001 samples per channel; both fault-related runs have 140,001.
The baseline fault-current peak is approximately $1.26\times10^{-6}$ kA,
while the phase-A peak for the physical-fault runs is approximately 1.552 kA.
The command-edit validation moves isolation from 5.0 to 5.1 s while retaining
the 6.0-s restoration time and the same fixed feeder model.

Together with the parameter manifests and fresh-output checks, this evidence
supports the implemented dashboard-to-PSCAD batch workflow. It is not a
measurement of end-to-end dashboard latency, proof of uninterrupted control
connectivity, or a live feedback demonstration. The project context also
records 114 passing automated tests; those tests were not rerun for this
read-only manuscript task and are not counted as additional experiments.
